\section{Beispiel Abschnitt}
\parencite{latex_companion} ist eine Beispielreferenz.
Es können Abkürzungen verwendet werden, wie \ac{BA} aber auch \ac{ABA}.

\subsection{Beispiel Unterabschnitt}

\begin{figure}[h]
    \centering

    \begin{subfigure}{.4\textwidth}
        \centering
        \includegraphics[width=.8\linewidth]{NAK_logo_full_text}
        \caption{Logo 1}
        \label{fig:subfig1}
    \end{subfigure}
    \begin{subfigure}{.4\textwidth}
        \centering
        \includegraphics[width=.8\linewidth]{NAK_logo_full_text}
        \caption{Logo 2}
        \label{fig:subfig2}
    \end{subfigure}

    \caption{Nordakademie Logos}
    \label{fig:fig}
\end{figure}

In Abb.~\ref{fig:subfig1} ist ein Beispiel Bild zu sehen.

\begin{listing}[H]
    \begin{minted}{c}
        int main() {
            // print hello world
            printf("hello, world!");
            return 0;
        }
    \end{minted}
    \caption{Hello world in C}
    \label{lst:listing}
\end{listing}

Zu Quelltext~\ref{lst:listing} kann auch Code innerhalb eines Textes ergänzt werden mit \mintinline{c}|int i = 1|.

\subsubsection{Beispiel Unterunterabschnitt}

\Blindtext{}
